\documentclass[10pt,a4paper]{amsart}
\usepackage[margin=19mm]{geometry}
\usepackage{amsmath,amssymb,mathtools}
\usepackage[scr=rsfs,cal=euler]{mathalfa}
\usepackage{xfrac}
\usepackage[unicode,hidelinks,bookmarks=false]{hyperref}
\newcommand{\ie}{i.e.\ }
\newcommand{\cf}{cf.\ }
\newcommand{\resp}{respectively\ }
\newcommand{\Subsection}{\subsection}
\providecommand{\nobreakdash}{\nobreak\hbox{-}}
\setlength{\emergencystretch}{2em}
\multlinegap0pt
\usepackage{bm}
\usepackage{bbm}
\usepackage{dsfont}
\usepackage{xspace}
\usepackage{cancel}
\usepackage{mathtools}
\usepackage{amsthm}
\makeatletter
\@ifundefined{theorem}{\newtheorem{theorem}{Theorem}}{}
\@ifundefined{lemma}{\newtheorem{lemma}[theorem]{Lemma}}{}
\@ifundefined{proposition}{\newtheorem{proposition}[theorem]{Proposition}}{}
\makeatother
\newcommand{\R}{{\mathbb R}}
\title[Spectrality of factors of product spectral measures]{Spectrality of factors of product spectral measures}
\author{Mihail N. Kolountzakis, Chun-Kit Lai, Kailing Lai, Jinjun Li}
\date{Source: arXiv:2607.01500v1 (CC BY 4.0)}
\begin{document}
\maketitle

\noindent\emph{Source and licence.} Spectrality of factors of product spectral measures. Version arXiv:2607.01500v1 (CC BY 4.0), distributed under the Creative Commons Attribution 4.0 International licence, as recorded in the arXiv licence metadata for this exact version and shown on the abstract page https://arxiv.org/abs/2607.01500v1. Full rights notice: licenses/arxiv-2607.01500-CC-BY-4.0.txt.\par\smallskip

\noindent\textit{Editorial scope.} Counted unit: the Proposition whose source label is \texttt{mapspectrum} in the article (source0.tex lines 486--489, source label \texttt{mapspectrum}), delivered together with the article's own complete original proof of it (source0.tex lines 490--523, one \texttt{proof} environment of 34 lines). No mathematics is written, repaired or re-derived: statement and proof are the article's own text line for line. Required context retained uncounted: measure (lines 395--397); partition (lines 402--404); 3or (lines 407--409); zero (lines 470--472); mapping (lines 477--483). Every definition, display and label of those lines is retained unchanged. Omitted from this package and not redistributed: 1--394, 398--401, 405--406, 410--469, 473--476, 484--485, 524--1051. Nothing inside a retained excerpt is omitted or altered: every retained body is delivered exactly as the article's lines are written, so no omission receipt of any kind accompanies this package. Citations: the retained excerpts carry no citation command at all, so no bibliography entry of the article is transcribed and no reference is redistributed. Reference adaptation: none was needed and none was made; every reference inside the delivered unit resolves inside it, so no unresolved reference or citation is produced. Printed slips kept literal: no printed word, symbol, hypothesis or number of the counted statement or of its proof was altered. Layout and class: the article's own class is \texttt{amsart} with packages including amssymb, bm, mathrsfs, mathtools, booktabs, enumerate, tikz, geometry, graphics, verbatim, graphicx, csquotes, enumitem, thmtools; the wrapper is the lane's pinned \texttt{amsart} editorial wrapper at 10pt on A4, which restates the theorem-like environments the retained text needs and adds the article's own macro definitions that the retained text needs (\texttt{\textbackslash R}), with extra wrapper packages (bm, bbm, dsfont, xspace, cancel, mathtools). The delivered numbering is the unit's own. Assets: no retained span contains a figure environment, a graphics-inclusion command, a PDF-page inclusion, a table or a TikZ picture; no image file is copied into this package, the package declares no asset dependency and no image file is redistributed with it. Licence: version arXiv:2607.01500v1 of this article is distributed under the Creative Commons Attribution 4.0 International licence, as recorded in the arXiv licence metadata for this exact version and shown on the abstract page https://arxiv.org/abs/2607.01500v1. A full rights notice is delivered at licenses/arxiv-2607.01500-CC-BY-4.0.txt. Characters: every retained excerpt is pure ASCII, so the delivered engine, XeLaTeX with its default Latin Modern text font, reports no missing character.
\begin{equation}\label{measure}
\rho = \mu\times \nu    
\end{equation}

\begin{equation}\label{partition}
\Lambda:=\Lambda_{1}\cup \Lambda_{2},\ \ \ \Lambda^{+}_\gamma:=\Lambda_{1}\cup (\Lambda_{2}+\gamma),\ \ \ \Lambda^{-}_\gamma:=\Lambda_{1}\cup (\Lambda_{2}-\gamma).
\end{equation}

\begin{lemma}\label{3or}Let $\Lambda$, $\Lambda^{+}_\gamma$ and $\Lambda^{-}_\gamma$ be the sets defined in $\eqref{partition}$, $\rho=\mu\times\nu$ be the measure defined in $\eqref{measure}$.
Suppose that $E(\Lambda)$ is an orthogonal basis for $L^2(\rho)$, and both $E(\Lambda^{+}_\gamma)$ and $E(\Lambda^{-}_\gamma)$ are orthogonal sets of $L^2(\rho)$, then each one of the sets $E(\Lambda^{+}_\gamma)$ and $E(\Lambda^{-}_\gamma)$ is an orthogonal basis for $L^2(\rho)$.
\end{lemma}

\begin{lemma}\label{zero}
Let $\lambda =(\lambda_1,\lambda_2)$ and $\lambda'=(\lambda'_1,\lambda_2')$ be two distinct points in $\mathbb{R}^{n+m}$, and $\rho=\mu\times\nu$ be the measure defined in \eqref{measure}. The exponential functions $e_{\lambda}$ and $e_{\lambda'}$ are orthogonal in $L^{2}(\rho)$ if and only if $\lambda_{1}-\lambda'_{1}\in \mathcal{Z}(\widehat{\mu})$, or $\lambda_{2}-\lambda'_{2}\in \mathcal{Z}(\widehat{\nu})$. %In particular, if $\mathcal{Z}(\hat{\mu})=\emptyset$, then the exponential functions $e_{\lambda}$ and $e_{\lambda'}$ are orthogonal in $L^{2}(\rho)$ if and only if $\lambda'_{2}-\lambda_{2}\in \mathcal{Z}(\hat{\nu})$.
\end{lemma}

\begin{equation}\label{mapping}
\phi_{t}(\lambda):=
\left\{\begin{matrix}
\lambda,\ \ \ \ \ \ \ \ \ & \ \lambda_{1}\in  \Gamma,\\
\lambda+\tau(t),& \lambda_{1}\notin \Gamma,
\end{matrix}\right.
\end{equation}

\begin{proposition}\label{mapspectrum}
 Given $t\in \R^n$. Let $\phi_{t}$ be the mapping defined by $\eqref{mapping}$, $\rho=\mu\times\nu$ be the measure defined in $\eqref{measure}$, and ${\mathcal Z}(\widehat{\mu})\cup \{0\}$ is a subset of  the {additive} subgroup $\Gamma$. If $\Lambda$ is a spectrum of $\rho$, then
$\phi_{t}(\Lambda)$ is also a spectrum of $\rho$.
\end{proposition}

\begin{proof}
%It is obvious that the mapping $\phi_{\gamma}$ from $\Lambda$ to $\phi_{\gamma}(\Lambda)$ is surjective. The proof considers two cases, $\mathcal{Z}(\hat{\mu})\neq\emptyset$ and $\mathcal{Z}(\hat{\mu})=\emptyset$. Since the proof for case $\mathcal{Z}(\hat{\mu})=\emptyset$ is analogous to that for case $\mathcal{Z}(\hat{\mu})\neq\emptyset$, we only present the proof for case $\mathcal{Z}(\hat{\mu})\neq\emptyset$.

We first prove that $E(\phi_t(\Lambda))$ is mutually orthogonal in $L^2(\rho)$. Given distinct $\lambda = (\lambda_1,\lambda_2)$ and $\lambda' = (\lambda_1',\lambda_2')$ in $\Lambda$. We have three cases to consider according to the definition of $\phi_t$. 
\begin{enumerate}
    \item  $\lambda_1,\lambda_1'\in\Gamma$;
    \item  $\lambda_1,\lambda_1'\not\in\Gamma$;
    \item  $\lambda_1\in\Gamma$ and $\lambda_1'\not\in \Gamma$ or vice versa. 
\end{enumerate}

In the first two cases, $\phi_t(\lambda)-\phi_t(\lambda') = \lambda-\lambda'$. As $e_{\lambda}$, $e_{\lambda'}$ are mutually orthogonal in $L^2(\rho)$, we know the same is also true for $e_{\phi_t(\lambda)}$, $e_{\phi_t(\lambda')}$.

In the last case, we only need to consider one situation since the role of $\lambda_1,\lambda_1'$ is symmetric.  Now,
\begin{equation}\label{phi_t-difference}
\phi_t(\lambda)-\phi_t(\lambda') = (\lambda_1-\lambda_1'-t,\lambda_2-\lambda_2').
\end{equation}
From our assumption, $\lambda_1\in \Gamma$, while $\lambda_1'\not\in \Gamma$. As  $\Gamma$ is an {additive} group, we have $\lambda_1-\lambda_1'\not\in\Gamma$. By Lemma \ref{zero} and ${\mathcal Z}(\widehat{\mu})\subset \Gamma$, it means that $\lambda_1-\lambda_1'\not\in{\mathcal Z}(\widehat{\mu})$ and thus $\lambda_2-\lambda_2'\in{\mathcal Z}(\widehat{\nu})$. Hence, $\phi_t(\lambda)-\phi_t(\lambda')\in {\mathcal Z}(\widehat{\rho})$ in view of the second component in (\ref{phi_t-difference}). This completes the proof of the mutual orthogonality. A similar argument also  shows that  $E(\phi_{-t}(\Lambda))$ is also a mutually orthogonal set. 


 We next show that $\phi_{t}(\Lambda)$ is indeed a spectrum of $\rho$. To this end, we  divide the set $\Lambda$ into two subsets, $\Lambda_{1}$ and $\Lambda_{2}$, such that they satisfy $\eqref{partition}$ and the assumptions in Lemma \ref{3or}.

Let
\[
\Lambda_{1}:=\{\lambda\in \Lambda:\lambda_{1}\in \Gamma\},\ \ \Lambda_{2}:=\{\lambda\in \Lambda:\lambda_{1}\notin \Gamma\}.
\]
We have immediately that 
$$
\phi_{t}(\Lambda)=\Lambda_{1}\cup(\Lambda_{2}+\tau(t)), \ \ \phi_{-t}(\Lambda)=\Lambda_{1}\cup(\Lambda_{2}-\tau(t)).$$
Moreover,
$$\Lambda_{1}\cap(\Lambda_{2}+\tau(t))=\varnothing, \ \ \Lambda_{1}\cap (\Lambda_{2}-\tau(t))=\varnothing.$$
Indeed, if either one of the the intersections is not disjoint, we will find some $(\lambda_1,\lambda_2) = (\lambda_1'+t,\lambda'_2)$ or $(\lambda_1'-t,\lambda'_2)$ where $(\lambda_1',\lambda_2')\in\Lambda_2$. In both cases, it implies that $\lambda_2 = \lambda_2'$. Since $(\lambda_1-\lambda_1',\lambda_2-\lambda_2')\in{\mathcal Z}(\widehat{\rho})$,   $\lambda_1-\lambda_1'\in{\mathcal Z}(\widehat{\mu})\subset \Gamma$ by the mutually orthogonality of $E(\Lambda)$ and Lemma \ref{zero}. As $\lambda_1\in\Gamma$, this forces $\lambda_1'\in\Gamma$ which is a contradiction since $\lambda_1'\in\Lambda_2$.

Finally, we can invoke  Lemma \ref{3or} to conclude that {each of} $\phi_{t}(\Lambda)$ and $\phi_{-t}(\Lambda)$ is a spectrum of $\rho$.
\end{proof}

\end{document}
